From a conventional neuroscience perspective, the body is often described in terms of discrete events: spikes, synaptic transmissions, and localized chemical signals. That view has been enormously productive, but it can also obscure a complementary picture in which living systems behave as coupled, multiscale oscillators. In this framing, physiology is not only a sequence of isolated outputs; it is also a continuous pattern of timing, phase, amplitude, and coherence. The body can therefore be understood as a resonant system in which structure and function are organized by recurring rhythms.

This shift from spike-centric neuroscience to field-aware models does not reject the importance of action potentials or molecular signaling. Rather, it places them within a broader dynamical context. Membranes exhibit electrical oscillations, mitochondria maintain rhythmic metabolic activity, and neural circuits generate coordinated rhythms across local and distributed networks. These oscillations occur at different scales, yet they are not independent. They interact through feedback, coupling, and entrainment, producing emergent patterns that can stabilize or destabilize physiological states.

At the cellular level, membranes are not merely passive barriers. Their electrical properties support oscillatory behavior that influences excitability, transport, and signaling. Mitochondria likewise participate in rhythmic processes tied to energy production, redox balance, and intracellular coordination. At the systems level, neural rhythms organize perception, attention, sleep, and motor control. Taken together, these examples suggest that biological function is deeply shaped by resonance across scale, from subcellular dynamics to whole-brain activity.

This perspective also invites attention to pattern formation and stability. In physical systems, resonant interactions can produce motifs that persist as attractors within a dynamic landscape. A similar logic may apply in biology, where recurring spatial and temporal patterns can become stabilized through repeated coupling. Cymatics offers a useful analogy here: when a medium is driven at particular frequencies, it can organize into visible geometric forms. While living tissue is far more complex than a vibrating plate or fluid surface, the analogy highlights an important principle: oscillatory input can shape structure, and structure can in turn constrain oscillatory behavior.

In this sense, intracellular field structure may be understood not as a static architecture but as a dynamic patterning process. Stable motifs can emerge when multiple oscillators lock into compatible relationships, creating attractor-like states that support function. This does not imply that biology is reducible to simple harmonic motion. Rather, it suggests that resonance is one of the organizing principles through which living systems maintain coherence, adapt to perturbation, and transition between states.

The sections that follow develop this idea into a practical framework. If the body is a resonant system, then targeted frequencies may be able to influence specific physiological processes by engaging the oscillatory structures already present in membranes, organelles, tissues, and neural networks.
